\begin{table*}[t]
\caption{Matched encoder comparison on screened CFP. Rates are percentages; ASR differences are percentage points.}
\label{tab:encoder-cfp}\centering\footnotesize
\begin{tabular}{@{}rlrrrrr@{}}\toprule
Bits & Method & MBF TAR / FAR & R50 TAR / FAR & MBF ASR & R50 ASR & ASR difference [95\% CI]\\\midrule
128 & RO & 39.33 / 0.134 & 65.50 / 0.133 & 19.66 & 38.46 & 18.80 [16.68, 20.86]\\
128 & SRHT & 37.76 / 0.113 & 65.12 / 0.120 & 18.84 & 38.28 & 19.43 [17.03, 21.81]\\
128 & PCA-guided RO & 43.50 / 0.115 & 66.94 / 0.103 & 32.96 & 53.09 & 20.13 [17.74, 22.68]\\
512 & RO & 82.52 / 0.172 & 94.61 / 0.134 & 78.02 & 92.30 & 14.28 [11.13, 17.52]\\
512 & SRHT & 81.95 / 0.167 & 94.67 / 0.128 & 77.52 & 92.23 & 14.71 [11.67, 17.95]\\
512 & PCA-guided RO & 82.38 / 0.200 & 94.37 / 0.130 & 78.56 & 92.03 & 13.48 [10.34, 16.72]\\
\bottomrule\end{tabular}
\tabnote{MBF: MobileFaceNet; R50: ResNet50 ArcFace. Each encoder has a separate validation-only threshold targeting 0.1\% FAR. ASR is calibrated cross-image fresh-key acceptance at $K=10$. The fresh operator is hidden from reconstruction. Paired percentile intervals use 2000 identity resamples and condition on the fitted pipelines; they are pointwise and not multiplicity-adjusted. These replication runs remain separate from the original baseline tables.}
\end{table*}
